\documentclass[]{article}
\usepackage[dvips]{graphicx}
\usepackage[T1]{fontenc}
\begin{document}

\title{The \LaTeX{} DTD and multiple languages}
\author{Michel Goossens
\thanks{Partly from an example in \emph{The \LaTeX{} Companion}}}

\date{August 4th, 1998}

\maketitle
\section{The basic principles\label{sec-en}}

\par

This is an example input file. We start in English to show the
principle. You should especially pay attention that we have used
slightly different notation for some of the common \LaTeX{} constructs,
such as the dashes, which come in three sizes: an intra-word dash, a
medium dash for number ranges like 1--2, and a punctuation
dash---like this.  Text can be emphasized as \emph{shown
here}. An ellipsis is made with \dots Footnotes\footnote{This
is a simple footnote.
\par
It can also contain \texttt{par}
elements.} are tricky constructs, since one must be
careful not to nest them.
\subsection{Dealing with special characters}

\par

XML has a different set of reserved characters than \LaTeX{}, in
particular, when you want to use any of the three characters
\texttt{\&}, \texttt{<}, 
and \texttt{>}, you should enter them as
\texttt{\&amp;}, \texttt{\&lt;}, and
\texttt{\&gt;}, respectively.
\subsection{\LaTeX{} and mathematical formulae\label{sec-math}}

\par

\LaTeX{} and \emph{a fortiori} \TeX{} are very good at typesetting
mathematical formulae, like 
$x- 3 y + z < 7$ or 
$a_{1} > x^{2n} + y^{2n} > x'$ or 
$(A, B) = \sum_{i} a_{i} b_{i}$.
Do not forget that for reasons of consistency, if you want to refer to
a variable in one of the formulae, such as the symbol 
$x$, you must also use math mode in the text.
\section{Beispiel eines Textes in deutscher Sprache}
\subsection{Eine EPS Abbildung}

\par

Dieser Abschnitt zeigt, wie man eine PostScript-Abbildung
\cite{bib-PS} in ein Dokument einbinden kann.  
Abbildung~\ref{fig-psfig} wurde mit dem Befehl 

\begin{verbatim}
\includegraphics[width="3cm"]{file="colorcir.eps}
\end{verbatim}
in den Text aufgenommen.

\begin{figure}\centering

\includegraphics[width=3cm, scale=.5]{colorcir.eps}
\caption{Ein EPS Bild\label{fig-psfig}}
\end{figure}
\subsection{Beispiel einer Tabelle}

\par
Die Tabelle~\ref{tab-exag} auf Seite~\pageref{tab-exag} zeigt eine Tabelle.

\begin{table}\centering
\caption{
Eingabe der deutschen Zusatzzeichen in \LaTeX{}\label{tab-exag}}

\begin{tabular}{ccccccc}
\texttt{"a}~�&\texttt{"A}~�&\texttt{"o}~�&\texttt{"O}~�&\texttt{"u}~�&\texttt{"U}~�&\texttt{"s}~�\\

\end{tabular}
\end{table}
\section{Continuation du texte en fran�ais}
\subsection{Traiter les listes\label{sec-list}}

\par

Les listes sont utilis�es fr�quemment pour structurer ou mettre
en �vidence certains �l�ments d'un document (voir \cite{bib-Liste}).

\begin{itemize}

\item Ceci est le premier �l�ment d'une liste non-ordonn�e. Chaque �l�ment
    de ce type de liste est pr�c�d� d'un signe distinctif, comme une
    puce, un tiret, etc.
  
\item Ce second �l�ment de la m�me liste contient une liste de
    \emph{description} imbriqu�e.
    
\begin{description}

\item[XML]Meta langage pour d�finir des classes de documents
\item[XLL]Langage pour d�finir des hyperliens entre diff�rentes
      parties de documents XML
\end{description}

    Nous continuons notre texte � l'int�rieur de la premi�re liste.
  \end{itemize}

\begin{thebibliography}{99}
\bibitem{bib-PS}
Adobe Inc. \emph{PostScript Handbuch (2. Auflage)}
           Addison-Wesley (Deutschland) GmbH, Bonn, 1991.
\bibitem{bib-Liste}
Michel Goossens. Personnaliser les listes \LaTeX{}.
           \emph{Cahiers GUTenberg}, 17:32--48, mai 1994.

\end{thebibliography}
\end{document}
